\section{Expanding Binomial Expressions}

Hopefully you recall the work we did in Year 1 on the binomial expansion. We
used the Pascal triangle and the Combination method to expand and simplify
certain expressions. Let's remind you of this in the activity below.

\begin{activity}{Activity 1.2: Revision}
\begin{enumerate}
  \item In groups or pairs, write down or generate the Pascal triangle up to $(x+y)^5$.
  \item Use your results in step 1 to simplify the expressions:
  \begin{enumerate}
    \item[a.] $(2x+y)^3$
    \item[b.] $(a-3b)^4$
    \item[c.] $(2+\sqrt{x})^5$
  \end{enumerate}
  \item Apply the combination method to also simplify the above tasks.
  \item Compare your results in (2) and (3) and write down any observations.
  \item Show your solutions, or write up, to classmates in other groups or to
  your mathematics teacher.
\end{enumerate}
\end{activity}

The Binomial Expansion helps us to expand and simplify complex expressions,
which would have been difficult using the idea of algebraic expressions. Under
algebraic expressions we could only expand with a positive integer power, but the
binomial theorem will help us to expand any expression with a rational number
as the power.

\begin{activity}{Activity 1.3 -- Derivation of Binomial Theorem}
\begin{enumerate}
  \item Look for the $^{n}C_{r}$ button on your calculator.
  \item Task: use $^{n}C_{r}$ on your calculator to evaluate
  \begin{enumerate}
    \item[i)] $^{3}C_{0}$
    \item[ii)] $^{3}C_{1}$
    \item[iii)] $^{3}C_{2}$
    \item[iv)] $^{3}C_{3}$
  \end{enumerate}
  \item Hopefully the answers you have are: $1\ 3\ 3\ 1$. This was used in Year 1 to
  expand certain expressions.
\end{enumerate}
\end{activity}

The combination of $n$ items taking $r$ at a time is:
\[
  \binom{n}{r} = {}^{n}C_{r} = \frac{n!}{(n-r)!\,r!}, \qquad \text{where } n! \text{ is read as ``}n\text{ factorial''}
\]
and is explained as $n! = n(n-1)(n-2)(n-3)\cdots 2 \times 1$.

For example, $5! = 5(5-1)(5-2)(5-3)(5-4) = 5 \times 4 \times 3 \times 2 \times 1 = 120$.

You could also find $5!$ using a calculator. Look for $!$ on your calculator to
find $5!$, to confirm the answer obtained above.

\begin{example}{Example 1.9}
In groups or in pairs, solve the following without using a calculator:
\begin{enumerate}
  \item[a.] $^{6}C_{3}$
  \item[b.] $^{4}C_{2}$
\end{enumerate}
Confirm your answers with a calculator.

\medskip\noindent\textbf{Solution}
\begin{enumerate}
  \item[a.] $^{6}C_{3} = \dfrac{6!}{(6-3)!\,3!} = \dfrac{6\times5\times4\times3\times2\times1}{3\times2\times1\times3\times2\times1} = \dfrac{6\times5\times4}{3\times2\times1} = \dfrac{120}{6} = 20$
  \item[b.] $^{4}C_{2} = \dfrac{4!}{(4-2)!\,2!} = \dfrac{4\times3\times2\times1}{2\times1\times2\times1} = \dfrac{4\times3}{2\times1} = \dfrac{12}{2} = 6$
\end{enumerate}
\end{example}
